\documentclass[12pt]{article}

\usepackage{amssymb,amsmath,amsthm} % Include AMS Math symbols
\usepackage{graphicx}               % Allows for importation of figures 

\usepackage{booktabs} % professional table typesetting commands 

\usepackage{color}   % Sets up a nice typesetting environment for Python
\usepackage{listings}
\definecolor{dkgreen}{rgb}{0,0.6,0}
\definecolor{gray}{rgb}{0.5,0.5,0.5}
\definecolor{mauve}{rgb}{0.58,0,0.82}
\lstset{frame=tb,
  language=Python,
  aboveskip=3mm,
  belowskip=3mm,
  showstringspaces=false,
  columns=flexible,
  basicstyle={\small\ttfamily},
  numbers=left,
  numberstyle=\tiny\color{gray},
  keywordstyle=\color{blue},
  commentstyle=\color{dkgreen},
  stringstyle=\color{mauve},
  breaklines=true,
  breakatwhitespace=true,
  tabsize=3
}

% define command \expandableinput  for table input
\makeatletter
\def\expinput#1{\@@input#1 }
\makeatother

\title{ HW \# \\ Math/CS 471, Fall 202*}
\author{Student 1 Full Name \\ Student 2 Full Name}
\date{\today}   % Activate to display a given date or no date


\begin{document}

% Display title
\maketitle
\clearpage

\section{Question 1: First homework question of the year, very excited}

\subsection{What}
Write briefly what the question is (what you are asked to do).  Someone with a similar background to you should be able to read this section and understand what the problem is.  Do not assume that the reader already knows what the problem is.  This is part of writing scientific reports, to provide some introduction to what you are doing.

\subsection{How / Method / Approach}
Describe how you solved the question, and show your results (figures, tables, numbers, etc) 

\subsection{Discussion / Interpretation} 
Discuss your results.  Did the method perform as expected?  Did you change parameters and 
investigate some kind of algorithm or problem specific behavior?  What did you observe and
can you explain it?

This does not need to always be long, but some discussion or interpretation is needed.

\section{Question 2: Question Title}

Repeat same report structure for all following questions\\\\


\section*{Useful Latex Examples To Get You Started}

\subsection{Equations}
It has very powerful support for equations such as these:
\begin{equation}
f(x) = e^{1/2 - \sin(5 \pi    x)}, \label{eq:overcomplicated_exp_function} 
\end{equation}
\begin{equation}
p(x) = 1+x+x^2. \label{eq:special_polynomial} 
\end{equation}
% The label command is used to provide a unique identifier for each equation that is independent of any numbering

\noindent
Here is how to typeset fractions and square-roots
\begin{equation}
    r = \frac{-b + \sqrt{b^2 - 4ac}}{2a}
\end{equation}

\noindent
For inline equation, you would use the dollar sign, i.e., $p(x) = 1 + x + x^2$.

\noindent
If you want an un-numbered equation use the asterisk
\begin{equation*}
    q(x) = x^3
\end{equation*}

\noindent
If you want a multi-line equation use ampersands for the alignment column and then two backslashes to denote a line break.
\begin{align}
   p(x) & = (x + 1)(x + 2)(x) \\
        & = (x^2 + 3x + 2)(x) \\
        & = x^3 + 3x^2 + 2x
\end{align}

\subsection{Referencing}
\LaTeX also supports a very convenient labeling and referencing system
through the command \texttt{\\label} and \texttt{\\ref}.  Above, equation
(\ref{eq:special_polynomial}) is a polynomial and the function described by equation
(\ref{eq:overcomplicated_exp_function}) is plotted in Figure \ref{fig:1}. 

{\bf Note} that you may have to compile the document two times to get
all the cross-referencing between labels to work correctly.

\subsection{Figures}
Figures are also easy to incorporate. It is often best to use a figure
format that is based on vector graphics, e.g. Encapsulated PostScript
(eps). PNG is another good format.  An example is given in Figure \ref{fig:1}. 

\begin{figure}[htb]
\begin{center}
\includegraphics[width=0.75\textwidth]{sample.png}
   \caption{This figure displays the function from equation (\ref{eq:overcomplicated_exp_function}).}
    \label{fig:1}
\end{center}
\end{figure}

\subsection{Tables}
Tables are often useful to display data more compactly than figures.

\begin{table}[h]
\begin{centering}
% We have two total columns in the table
% The "pipe" characters "|"  denote a vertical line will divide the columns
%                  left-aligned column   center-aligned column
  \begin{tabular}{      l             |        c             }
   \toprule  % horizontal line at top of table 
   \expinput{meals.tex}  
   \bottomrule  % horizontal line at bottom of table 
  \end{tabular}
  \caption{A table of delicious meals. \label{tab:1}}
  \end{centering}
\end{table}


\subsection{Code Typesetting}

Simple typesetting
\begin{verbatim}
from scipy import *
from matplotlib import pyplot as plt
a = 4
b = 2
c = a*b
if c > 2:
    x = linspace(0,1,100)
    y = sin(x)
    plt.plot(x,y)
    plt.show()
\end{verbatim}


Fancy typesetting for code
\begin{lstlisting}
from scipy import *
from matplotlib import pyplot as plt
a = 4
b = 2
c = a*b
if c > 2:
    x = linspace(0,1,100)
    y = sin(x)
    plt.plot(x,y)
    plt.show()
\end{lstlisting}

%\bibliographystyle{plain}
%\bibliography{bibfile}

\end{document}  

